\section{Gyllir}
\label{sec:basics:gyllir}

Calling \token{gyc} by hand stops being reasonable at about the second source
file. Gyllir is the build system and package manager that takes over from there:
it assembles the command lines, tracks the version of the project, and fetches
the libraries it depends on. Nothing in Ymir requires it, and this book uses
whichever of the two tools is the more convenient -- \token{gyc} directly for the
short single-file listings that make up most of the examples, and whenever the
compiler's intermediate representations are on show, as in the previous section;
Gyllir as soon as a program spans several files or reaches for an external
library.

A new project starts with the \textit{init} subcommand, which asks a handful of
questions, each with a default, and writes out an empty project around the
answers.

\begin{lstlisting}[style=bashVerb, escapechar=@+]
@+\textcolor{teal!80}{alice@dev:\mtilde}@+$ mkdir my-project
@+\textcolor{teal!80}{alice@dev:\mtilde}@+$ cd my-project
@+\textcolor{teal!80}{alice@dev:\mtilde}@+@+\textcolor{green!60!black}{/my-project}@+$ gyllir init
Name [main]: my-project
Author name [Alice Doe <alice-doe@mail.com>]:
Description [A minimal Ymir app]:
License [proprietary]:
Type (executable/library) [executable]:
Registry [local:/home/alice/.local/gyllir/my-project]:
\end{lstlisting}

That leaves three files in the current directory.
\textit{gyllir.toml} is the configuration Gyllir reads on every invocation; with
the defaults above it describes an executable. \textit{src} holds the source
files -- for now a single \textit{main.yr} with placeholder content -- and
\textit{test} holds the unit tests, the subject of a later chapter.

\begin{lstlisting}[style=bashVerb, escapechar=@+]
@+\textcolor{teal!80}{alice@dev:\mtilde}@+@+\textcolor{green!60!black}{/my-project}@+$ tree
@+\color{teal}\hspace{0.5ex}\raisebox{0.5ex}{\texttt{.}}\vspace{-3pt}@+
@+\japan{├──} \raisebox{0.5ex}{\texttt{gyllir.toml}}\vspace{-3.5pt}@+
@+\japan{├──}~\color{teal}{\raisebox{0.5ex}{\texttt{src}}}\vspace{-3.5pt}@+
@+\japan{│~~~~~└──}~\raisebox{0.5ex}{\texttt{main.yr}}\vspace{-3.5pt}@+
@+\japan{└──}~\color{teal}{\raisebox{0.5ex}{\texttt{test}}}\vspace{-3.5pt}@+
@+\japan{~~~~~~~└──}~\raisebox{0.5ex}{\texttt{\_\_test\_\_.yr}}\vspace{-3.5pt}@+

2 directories, 3 files
\end{lstlisting}

\textit{gyllir build} produces the executable; \textit{gyllir run} produces it
and then runs it.

\begin{lstlisting}[style=bashVerb, escapechar=@+]
@+\textcolor{teal!80}{alice@dev:\mtilde}@+@+\textcolor{green!60!black}{/my-project}@+$ gyllir run
  @+\color{black!30!applegreen}{Compiling}@+ my-project v0.1.0
  @+\color{black!30!applegreen}{Compiling}@+ 1 modules in thread 1 (among 1 changed files of 1)
  @+\color{black!30!applegreen}{Produced}@+ executable file my-project
  @+\color{black!30!applegreen}{Executing}@+ project my-project v0.1.0

Hello World!
\end{lstlisting}

\vfill
\pagebreak
